\clearpage
\phantomsection
\addcontentsline{toc}{chapter}{Abstract}
\thispagestyle{plain}

\vspace*{0.2cm}

\begingroup
\singlespacing
\noindent
\begin{minipage}[t]{0.45\textwidth}
\raggedright
{\small Thesis advisor\par}
\vspace{0.12cm}
{\normalsize\bfseries \AdvisorName\par}
\end{minipage}
\hfill
\begin{minipage}[t]{0.45\textwidth}
\raggedleft
{\small Author\par}
\vspace{0.12cm}
{\normalsize\bfseries \AuthorName\par}
\end{minipage}
\endgroup

\vspace{0.6cm}

\begin{center}
\begingroup
\singlespacing
{\normalsize\bfseries \ThesisTitle\par}
\vspace{0.45cm}
{\normalsize\bfseries Abstract}
\endgroup
\end{center}

\vspace{0.35cm}

Test code is essential for software quality, yet software maintenance research has focused mainly on production code, leaving important aspects of test-code quality and evolution underexplored. This work examines two dimensions of test-code maintenance: self-admitted technical debt (SATD) and method-level evolution.
To understand explicit maintenance concerns, comments from open-source Java projects were manually analyzed. SATD comments were identified and classified into categories, including test-specific ones. Existing SATD detection tools and both open-source and proprietary large language models were unreliable for detecting test-code SATD.
Studying test-code evolution requires accurate histories of test and production methods. Existing method-history generation tools, however, were evaluated against inaccurate oracles and mainly on production methods, leaving their reliability uncertain. To address this, updated and new oracles are constructed, and HistoryFinder is introduced as a method-history generation tool that achieves the highest overall accuracy with competitive runtime and outperforms other research-based tools. To examine test-production relationships, a test-to-production method-linking ground truth was constructed and used to evaluate existing linking techniques. Using these links and method histories, the analysis shows that some test methods are revised more frequently than their corresponding production methods and that several test-smell types are more prevalent among highly revision-prone test methods.
Overall, this work contributes ground truths, tools, and empirical evidence for understanding test-code quality and evolution.
